% !TEX root = chapter-standalone.tex

\section{Associative stacking categories}

\todotext{J: discuss associators and distinguish strict and non-strict case and talk about strictification... (maybe state the strictification theorem somewhere also, as a way to illustrate/motivate the construction we then use...)}

%\begin{ctdefinition}[Associative stacking category]
%    \label{def:assoc-stacking-cat}
%    An \maindef{associative stacking category} is a \SY{stacking category} $\tup{\CatC, \mtimescatob, \mtimescatmor}$ satisfying the following.
%
%    \condit
%    \begin{itemize}
%        \item The stacking operation~$\mtimescatob$ makes the objects into a \SY{semigroup}.
%        \item The stacking operation~$\mtimescatmor$ makes the morphisms into a \SY{semigroup}.
%    \end{itemize}
%\end{ctdefinition}

\begin{ctdefinition}[Strict associative stacking category]
    \label{def:strict-assoc-stacking-cat}
    An \maindef{strict associative stacking category} is

    \constit
    \begin{enumerate}
        \item a \SY{functorial stacking category} $\tup{\CatC, \mtimescat}$;
    \end{enumerate}

    \condit
    \begin{enumerate}
        \item the two composite functors $(( - ) \mtimescat ( - )) \mtimescat ( - )$ and $( - ) \mtimescat (( - ) \mtimescat ( - ))$ are equal as functors $\CatC \times \CatC \times \CatC \to \CatC$.
    \end{enumerate}
\end{ctdefinition}

\todotext{J: make commutative diagram version of the above associativity condition}

\begin{ctdefinition}[Associative stacking category]
    \label{def:assoc-stacking-cat}
    An \maindef{associative stacking category} is a \

    \constit
    \begin{enumerate}
        \item \SY{functorial stacking category} $\tup{\CatC, \mtimescat}$;
        \item a natural isomorphism
              \begin{equation}
                  \associator \colon (( - ) \mtimescat ( - )) \mtimescat ( - ) \nto ( - ) \mtimescat (( - ) \mtimescat ( - ))
              \end{equation}
    \end{enumerate}

    \condit
    \begin{enumerate}
        \item Pentagon identities: for all $\Obja,\Objb,\Objc,\Objd\setin \ObC$, the following diagram must commute
              \equationsag{30_pentagon_identity}{eq:30_pentagon_identity}
    \end{enumerate}
\end{ctdefinition}

If a \SY{semicategory} is \SY{associative stacking}, then the two ways of stacking three morphisms -- as depicted in \cref{fig:associative-stacking-1} and \cref{fig:associative-stacking-2}, respectively -- give the same result.
For \SY{associative stacking} \SY{semicategories} we will use a simpler diagrammatic notation, where the diagrams in \cref{fig:associative-stacking-1} and \cref{fig:associative-stacking-2} instead look like the ones in \cref{fig:associative-stacking-1-new} and \cref{fig:associative-stacking-2-new}.
Since these two diagrams depict that same morphism, when it is convenient we will also simply depict them as in \cref{fig:triple-stacked} (and similarly for any number of stacked morphisms).
In particular, in an \SY{associative stacking} \SY{semicategory}, and stacking of two morphisms will be depicted as in \cref{fig:double-stacked-new}.

\begin{marginfigure}[]
    \centering
    \includesag{parallelism_associative-stacking-1-new}
    %\includegraphics[scale=0.15]{parallelism_associative-stacking-1-new}
    \caption{Associative stacking of three morphisms, in one order.}
    \label{fig:associative-stacking-1-new}
\end{marginfigure}

\begin{marginfigure}
    \centering
    \includesag{parallelism_associative-stacking-2-new}
    %\includegraphics[scale=0.15]{parallelism_associative-stacking-2-new}
    \caption{Associative stacking of three morphisms, in another order.}
    \label{fig:associative-stacking-2-new}
\end{marginfigure}

\begin{marginfigure}
    \centering
    \includesag{parallelism-triple-stacked}
    %\includegraphics[scale=0.15]{parallelism_triple-stacked}
    \caption{Our string diagram notation for a triple stack.}
    \label{fig:triple-stacked}
\end{marginfigure}

\begin{figure}[h!]
    \centering
    \includesag{parallelism-associative-stacking-two-new}
    %\includegraphics[scale=0.15]{parallelism_associative-stacking-two-new}
    \caption{Associative stacking of two morphisms.}
    \label{fig:double-stacked-new}
\end{figure}

\begin{remark}
    In an \SY{associative stacking} \SY{semicategory} it follows that we can stack any number of morphisms without needing to bracket using parentheses:
    \begin{equation}
        \prfperiod{
            \moran{i} \colon \Objan{i} \mto \Objbn{i}, \qquad 1 \leq i \leq n,
        }{
            \mtimescatmor_{i=1}^n \moran{i} \colon \mtimescatob_{i=1}^n \Objan{i} \mto \mtimescatob_{i=1}^n \Objbn{i}.
        }
    \end{equation}
\end{remark}

\begin{example}\label{ex:assoc-stacking-semicat-integers-max}
    We let the integers $\wnumbers$ to be the set of objects of an associative stacking semicategory, and we say that there exists a unique morphism $\Obja \mto \Objb$ if and only if $\Obja \leq \Objb$.
    (We have already seen that this forms a semicategory; in fact, a category.) As our stacking operation for objects we set
    \begin{equation}
        \Obja \mtimescatob \Objb := \max(\Obja, \Objb)
    \end{equation}
    and for stacking morphisms we say that if $\mora \colon \Obja \mto \Objb$ and $\morb \colon \Objc \mto \Objd$ exist, then there exists a unique morphism
    \begin{equation}
        \mora \mtimescatmor \morb \colon \Obja \mtimescatob \Objc \mto \Objb \mtimescatob \Objd
    \end{equation}
    which corresponds to (and is consistent with) the inequality
    \begin{equation}
        \max(\Obja, \Objc) \leq \max(\Objb, \Objd).
    \end{equation}
\end{example}

\begin{example}\label{ex:assoc-stacking-semicat-lists-concat}
    Let $\setA$ be a non-empty set and consider a semicategory where the collection of objects is $ \listsof \setA$, the set of non-empty lists of elements of $\setA$.
    We can define a pre-order on $\listsof \setA$ by setting, for any lists $\Obja, \Objb$,
    \begin{equation}
        \Obja \leq \Objb \definedas \text{length}(\Obja) \leq \text{length}(\Objb).
    \end{equation}
    Then we view this pre-order as a semicategory, and define the following stacking operations.
    Given lists $\Obja, \Objb$, let $\Obja \mtimescatob \Objb$ be the concatenation of $\Obja$ and $\Objb$, and given morphisms $\mora \colon \Obja \mto \Objb$ and $\morb \colon \Objc \mto \Objd$ representing inequalities, we let $\mora \mtimescatmor \morb \colon \Obja \mtimescatob \Objc \mto \Objb \mtimescatob \Objd$ represent the inequality
    \begin{equation}
        \text{length}(\Obja \mtimescatob \Objc) \leq \text{length}(\Objb \mtimescatob \Objd).
    \end{equation}

\end{example}

\subsection{One-object associative stacking semicategories}

\todotext{Maybe remove this subsection? or move later when talk about functorial}
Consider a special kind of a stacking semicategory $\tup{\CatC, \mtimescatob, \mtimescatmor}$ where the semicategory $\CatC$ has only one object (call it $\Obja$, say).

In this special case, we have $\mtimescatob_{i=1}^n \Objan{i} = \Obja$ for any $n \setin \natnumbers$.
The only hom-set is $\Hom_\CatC(\Obja, \Obja)$ and this is equipped with the stacking operation
\begin{equation}
    \mtimescatmor \colon \Hom_\CatC(\Obja, \Obja) \cartprod \Hom_\CatC(\Obja, \Obja) \to \Hom_\CatC(\Obja, \Obja)
\end{equation}
which makes $\Hom_\CatC(\Obja, \Obja)$ into a semigroup.
In other words, this means that an associative stacking semicategory with one object may equivalently be described as set together with two operations -- serial composition and stacking -- that equip said set with two semigroup structures.

\begin{example}
    The integers $\wnumbers$, equipped with addition and multiplication as serial composition and stacking respectively, specify an associative stacking category with one object.
    Alternatively, we may also choose multiplication as our serial composition, and addition as our stacking.
\end{example}

\subsection{Design problems}

\begin{example}
    As long as the propulsion and life support systems below do not interact, we can simply tensor the two \SY{design problems} representing these systems into one, big co-design problem (\cref{fig:examplemonoidal}).
    \begin{figure}[h!]
        \centering
        \includesag{50_engine_tensor_1_2}
        \caption{Example of tensor of \SY{design problems}. }
        \label{fig:examplemonoidal}
    \end{figure}
\end{example}

